\subsubsection*{The upper limit in the slope integral}
\label{par:marked-height-limit}

The preceding rank estimate bounds the integrand in the slope comparison. To integrate it, we also need a cutoff independent of the chosen translate.

Let \(R\) be the direct row space at the refined level and write
\(\bar\kappa\) for the mean logarithm of the period scalar.
The target of the period map is \(R^\vee\).

\begin{proposition}[Uniform upper limit for the slope integral]
\label{prop:marked-frame-height}
For the selected family, use the integral models of
Section~\ref{subsec:division-level-pins} and the period shifts
of Section~\ref{sec:word-division}. There is a middle exponent \(c_{HH}\),
chosen as a function of the outer parameters and \(\epsilon\), for which
the mean \(s\) in Lemma~\ref{lem:frame-slope} satisfies
\begin{equation}\label{eq:marked-frame-height}
          s-\mu_{\min}(R^\vee)\le b x^{c_{HH}} .
\end{equation}
The bound holds at every point of the family above one threshold.
\end{proposition}

\begin{proof}
\proofstep{Apply the sharp row comparison before changing level.}
On the integral models of
Section~\ref{subsec:division-level-pins} the uncut numerator
bound of Section~\ref{subsec:division-length} holds for
every lift at the same point. Thus the projection and
functional comparisons in Section~\ref{sec:word-division}
use all their original place-to-tuple assignments. The
sharp primitive coordinate map takes the arithmetic lattice
onto \(R\) and carries its Gram at every infinite place to
the chosen Gram. It therefore preserves maximal slope:
it takes subspaces bijectively to subspaces and preserves
the covolume of the induced lattice on each one.

At each
embedding the period scalar of this comparison is the one
in the base-level formula. Its log mean is consequently
\(\bar\kappa\) in both. If \(C_f\ge1\) covers the base
functional estimate in Proposition~\ref{prop:descended-row-height},
level passage gives
\[
    \mu_{\max}(R)+\bar\kappa
       \le b x^{C_f}+14\log|G|\le b x^{C_f+1}
       \qquad (x\ge56).
\]
Here \(\log|G|\le4b\) and \((x-1)x^{C_f}\ge56\).
In this application the numerator is needed for the
surjectivity and functional bounds before passing
levels. \label{cor:frame-height-limit-1}

\proofstep{Reduce the upper limit to five pointwise terms.}
For each sequence \(u\) write \(M_u\) for its reference lattice,
\(N_u\) for its shifted lattice, and \(h_u\) for the numerical
height, including the period shift. The column rank bound and the ramified rank estimate give
\(\operatorname{rk}N_u>0\). Use \(\mu(N_u)\le\mu_{\max}(N_u)\)
in \eqref{eq:upper-slopes} and add and subtract \(\bar\kappa\).
For every \(u\) the result is
\[
 \begin{aligned}
 \mu(N_u)-h_u+\mu_{\max}(R)
 &\le \mu_{\max}(M_u)+(G_a+F-E_u)+CmB_0\\
 &\quad -(h_u+\bar\kappa)+(\mu_{\max}(R)+\bar\kappa).
 \end{aligned}
\]
Here \(C\) bounds the same error in every word of the selected
family. Each term on the right needs a common upper bound
before taking the rank-weighted mean.

\proofstep{Bound the two slope costs in terms of the rank parameters.}
Put \(P=1+f+k+m+t+N_m+L\) and \(M=\max(C_t,5)+3\),
where \(C_t\) is the outer rank exponent. The size conditions
give \(P\le x^M\) eventually. Indeed the terms other than
\(N_m\) are bounded by powers of \(x\) from the size bounds
on the host, the ranks and \(L=x^\epsilon\). Substitution of the chain rank into the sum gives the bound
\[
 P=1+f+k+m+t+\bigl(k+m(t+1)\bigr)+L\le 8x^{M-1}\le x^M
 \qquad(x\ge9) .
\]
Also \(B_0=fmbL\), \(a,b',S_c\le Pb\),
\(\log(kv)\le3P\), and \(\log T_u\le C_T' mL\)
by the count on the selected sequences.

In the reference maximal
slope estimate of Proposition~\ref{prop:level-zero-column-slopes}
the three upper costs are fixed multiples of
\[
 |T|a\log(kv),\quad
 f(b'+S_c)\big(k\log(kv)+\log T_u\big),\quad fk^2 L .
\]
Substitute \(|T|,f,k,m,L\le P\) in these three costs.
They are bounded, respectively, by
\[
 3P^3b,\qquad (6+2|C_T'|)P^4b,\qquad P^4 .
\]
For the middle term use
\(k\log(kv)+\log T_u\le(3+|C_T'|)P^2\).
The fixed linear-combination coefficients are absorbed at
large \(x\) using \(P\le x^M\), so the first term costs at
most \(b x^{4M+1}\). The third term \(CmB_0\) has the
same bound because \(mB_0=fm^2bL\le P^4b\).
\label{cor:frame-height-limit-2}

\proofstep{Bound the gain and the forced term before subtracting the centered exponent.}
The unweighted gain at \(p\le X''\) is at most
\(4m|T|a/p\) and the subtracted share is nonnegative.
Above \(X''\) there is no gain. Thus the squared-norm weight
gives
\[
 G_a\le 2\sum_{p\le X''}(\log p)\frac{4m|T|a}{p}
          \le16m|T|a\log(1+X'').
\]
The forced term of one lift is at most \(16P^7 b/p\) to that
cutoff. In fact its type and ramified negative contributions
are bounded above by zero.

On an inert orbit the coefficient
per label is bounded by \(2/p\) using
\(p/(p^2-1)\le2/p\). The sum has at most \(2f\) labels after
averaging around the orbit, each with size at most
\((k-1)4P^5 b\). This includes the shorter orbit with
coefficient \(1/(p^2-1)\). The zero cases contribute nothing.
For the forced term, the logarithmic prime weights give
\[
  \sum_{j=1}^{m}\sum_{p\le X''}(\log p)\frac{16P^7 b}{p}
  \ll m P^7 b\log(1+X'')\ll P^9 b .
\]
Here \(\log(1+X'')\le4P\) and \(m\le P\).
The full-support exponent estimate gives \(-E_u\le CmB_0\),
since each lift's top logarithm contribution is nonnegative.
The two bounds give the stated upper bound \(b x^{9M+2}\)
on the second term for sufficiently large \(x\). \label{cor:frame-height-limit-3}

\proofstep{Estimate the period-free height at each infinite place.}
The fourth term is the negative of the period-free height
in Section~\ref{sec:word-division}. Its three parts are
the negative top size term, the negative sum of kernel
terms, and the positive sum of lower hook terms. The first
costs \(O((1+|C_T'|+|C_{\rm sz}|)P^4b)\). Its sequence-dependent part has at most \(P\) terms of size \(2Pb\log T_u\); the size correction uses \(L_a\le4P\). The respective upper bounds are therefore
\[
 2P^2 b\log T_u,\qquad 4P^3 b L_a|C_{\rm sz}|\le16P^4b|C_{\rm sz}|.
\]
For a kernel at one real place write \(Y=11P^5b\).
The entries of the row weights obey \(|\nu_i|\le Y\)
by addition of the top, short and chain offsets.
Here \(\log Y\le16P\) and \(|\log\Omega|\le7P\).

For every \(\nu_i\), including those of modulus below one, the common lower estimate reads
\[
 \varphi(\nu_i)-\nu_i\log\Omega\ge -Y\log Y-Y-1-7PY .
\] There are at most \(P\) summands. At the same place
the positive hook numerator has at most \(P\) log gamma
terms with arguments in \([2,10P^5b]\).
The lower endpoint uses the gap above the last entry
\(-b_{j-1}\), and the upper endpoint uses the same weight
bounds. Each \(\log\Gamma(y)\le y\log y\) on this range.
The hook denominator consists of differences of the decreasing
integral weight shifted by the index difference. They are
at least one and can be dropped in an upper bound.

The two bounds at one place follow explicitly from \(Y=11P^5b\):
\[
 \begin{aligned}
 P(Y\log Y+Y+1+7PY)&=O(P^7b),\\
 P(10P^5b)\log(10P^5b)&=O(P^7b).
 \end{aligned}
\]
Here the first line bounds the negative kernel terms and the second
bounds the positive hook numerator. The kernel
corrections per lift are bounded by
\(|C_{\rm sz}|(f N_m^2 L_a+fm^3 N_m^2 L)\);
the hook correction is bounded by \(|C_{\rm sz}| f R^2
\log(b\,\operatorname{rd}(E)g_0g_1 R)\), with this
logarithm at most \(6P^2\). After summing over at most
\(P\) places and \(P\) lifts the cost is \(O(P^9b)\).
Consequently \(-(h_u+\bar\kappa)\le b x^{9M+1}\)
eventually. \label{cor:frame-height-limit-4}

\proofstep{Choose the exponent, then take the rank-weighted mean.}
Choose the middle height parameter \(c_{HH}\) at least
\[
     1+\max\{4M+1,9M+2,9M+1,C_f+1\}.
\]
The threshold for it uses only the outer parameters and
\(\epsilon\). Indeed the sharp functional exponent is
obtained from the base functional bound at the middle choice
made as a function of the outer parameters. The four
remaining powers just computed involve only \(C_t\).
Constants multiplying these powers, including the constants
of the upper slope and the size terms, enter the eventual
scale threshold. This permits later constants to depend
on the middle choice.

For \(x\ge9\) the sum of five terms
with exponents at most \(c_{HH}-1\) is at most \(b x^{c_{HH}}\).
With \(r_u=\operatorname{rk}N_u\), multiply by \(r_u\) and
divide the summed inequality by \(\sum_u r_u>0\). This gives
\[
     \frac{\sum_u r_u(\mu(N_u)-h_u)}{\sum_u r_u}
       +\mu_{\max}(R)\le b x^{c_{HH}}.
\]
The first term is \(s\), since the list in
Lemma~\ref{lem:frame-slope} includes slopes with multiplicities.
Hermitian duality gives \(\mu_{\min}(R^\vee)=-\mu_{\max}(R)\).
This gives the bound in the proposition, using the sharp
functional bound and the reference column slope at the same
selected point. \label{cor:frame-height-limit-5}
\end{proof}
